% Options for packages loaded elsewhere
% Options for packages loaded elsewhere
\PassOptionsToPackage{unicode}{hyperref}
\PassOptionsToPackage{hyphens}{url}
\PassOptionsToPackage{dvipsnames,svgnames,x11names}{xcolor}
%
\documentclass[
  10pt,
  letterpaper,
  DIV=11,
  numbers=noendperiod]{scrartcl}
\usepackage{xcolor}
\usepackage{amsmath,amssymb}
\setcounter{secnumdepth}{-\maxdimen} % remove section numbering
\usepackage{iftex}
\ifPDFTeX
  \usepackage[T1]{fontenc}
  \usepackage[utf8]{inputenc}
  \usepackage{textcomp} % provide euro and other symbols
\else % if luatex or xetex
  \usepackage{unicode-math} % this also loads fontspec
  \defaultfontfeatures{Scale=MatchLowercase}
  \defaultfontfeatures[\rmfamily]{Ligatures=TeX,Scale=1}
\fi
\usepackage{lmodern}
\ifPDFTeX\else
  % xetex/luatex font selection
  \setmainfont[]{TeX Gyre Pagella}
  \setmonofont[]{Fira Code}
  \setmathfont[]{TeX Gyre Pagella Math}
\fi
% Use upquote if available, for straight quotes in verbatim environments
\IfFileExists{upquote.sty}{\usepackage{upquote}}{}
\IfFileExists{microtype.sty}{% use microtype if available
  \usepackage[]{microtype}
  \UseMicrotypeSet[protrusion]{basicmath} % disable protrusion for tt fonts
}{}
\usepackage{setspace}
\makeatletter
\@ifundefined{KOMAClassName}{% if non-KOMA class
  \IfFileExists{parskip.sty}{%
    \usepackage{parskip}
  }{% else
    \setlength{\parindent}{0pt}
    \setlength{\parskip}{6pt plus 2pt minus 1pt}}
}{% if KOMA class
  \KOMAoptions{parskip=half}}
\makeatother
% Make \paragraph and \subparagraph free-standing
\makeatletter
\ifx\paragraph\undefined\else
  \let\oldparagraph\paragraph
  \renewcommand{\paragraph}{
    \@ifstar
      \xxxParagraphStar
      \xxxParagraphNoStar
  }
  \newcommand{\xxxParagraphStar}[1]{\oldparagraph*{#1}\mbox{}}
  \newcommand{\xxxParagraphNoStar}[1]{\oldparagraph{#1}\mbox{}}
\fi
\ifx\subparagraph\undefined\else
  \let\oldsubparagraph\subparagraph
  \renewcommand{\subparagraph}{
    \@ifstar
      \xxxSubParagraphStar
      \xxxSubParagraphNoStar
  }
  \newcommand{\xxxSubParagraphStar}[1]{\oldsubparagraph*{#1}\mbox{}}
  \newcommand{\xxxSubParagraphNoStar}[1]{\oldsubparagraph{#1}\mbox{}}
\fi
\makeatother

\usepackage{color}
\usepackage{fancyvrb}
\newcommand{\VerbBar}{|}
\newcommand{\VERB}{\Verb[commandchars=\\\{\}]}
\DefineVerbatimEnvironment{Highlighting}{Verbatim}{commandchars=\\\{\}}
% Add ',fontsize=\small' for more characters per line
\usepackage{framed}
\definecolor{shadecolor}{RGB}{241,243,245}
\newenvironment{Shaded}{\begin{snugshade}}{\end{snugshade}}
\newcommand{\AlertTok}[1]{\textcolor[rgb]{0.68,0.00,0.00}{#1}}
\newcommand{\AnnotationTok}[1]{\textcolor[rgb]{0.37,0.37,0.37}{#1}}
\newcommand{\AttributeTok}[1]{\textcolor[rgb]{0.40,0.45,0.13}{#1}}
\newcommand{\BaseNTok}[1]{\textcolor[rgb]{0.68,0.00,0.00}{#1}}
\newcommand{\BuiltInTok}[1]{\textcolor[rgb]{0.00,0.23,0.31}{#1}}
\newcommand{\CharTok}[1]{\textcolor[rgb]{0.13,0.47,0.30}{#1}}
\newcommand{\CommentTok}[1]{\textcolor[rgb]{0.37,0.37,0.37}{#1}}
\newcommand{\CommentVarTok}[1]{\textcolor[rgb]{0.37,0.37,0.37}{\textit{#1}}}
\newcommand{\ConstantTok}[1]{\textcolor[rgb]{0.56,0.35,0.01}{#1}}
\newcommand{\ControlFlowTok}[1]{\textcolor[rgb]{0.00,0.23,0.31}{\textbf{#1}}}
\newcommand{\DataTypeTok}[1]{\textcolor[rgb]{0.68,0.00,0.00}{#1}}
\newcommand{\DecValTok}[1]{\textcolor[rgb]{0.68,0.00,0.00}{#1}}
\newcommand{\DocumentationTok}[1]{\textcolor[rgb]{0.37,0.37,0.37}{\textit{#1}}}
\newcommand{\ErrorTok}[1]{\textcolor[rgb]{0.68,0.00,0.00}{#1}}
\newcommand{\ExtensionTok}[1]{\textcolor[rgb]{0.00,0.23,0.31}{#1}}
\newcommand{\FloatTok}[1]{\textcolor[rgb]{0.68,0.00,0.00}{#1}}
\newcommand{\FunctionTok}[1]{\textcolor[rgb]{0.28,0.35,0.67}{#1}}
\newcommand{\ImportTok}[1]{\textcolor[rgb]{0.00,0.46,0.62}{#1}}
\newcommand{\InformationTok}[1]{\textcolor[rgb]{0.37,0.37,0.37}{#1}}
\newcommand{\KeywordTok}[1]{\textcolor[rgb]{0.00,0.23,0.31}{\textbf{#1}}}
\newcommand{\NormalTok}[1]{\textcolor[rgb]{0.00,0.23,0.31}{#1}}
\newcommand{\OperatorTok}[1]{\textcolor[rgb]{0.37,0.37,0.37}{#1}}
\newcommand{\OtherTok}[1]{\textcolor[rgb]{0.00,0.23,0.31}{#1}}
\newcommand{\PreprocessorTok}[1]{\textcolor[rgb]{0.68,0.00,0.00}{#1}}
\newcommand{\RegionMarkerTok}[1]{\textcolor[rgb]{0.00,0.23,0.31}{#1}}
\newcommand{\SpecialCharTok}[1]{\textcolor[rgb]{0.37,0.37,0.37}{#1}}
\newcommand{\SpecialStringTok}[1]{\textcolor[rgb]{0.13,0.47,0.30}{#1}}
\newcommand{\StringTok}[1]{\textcolor[rgb]{0.13,0.47,0.30}{#1}}
\newcommand{\VariableTok}[1]{\textcolor[rgb]{0.07,0.07,0.07}{#1}}
\newcommand{\VerbatimStringTok}[1]{\textcolor[rgb]{0.13,0.47,0.30}{#1}}
\newcommand{\WarningTok}[1]{\textcolor[rgb]{0.37,0.37,0.37}{\textit{#1}}}

\usepackage{longtable,booktabs,array}
\usepackage{calc} % for calculating minipage widths
% Correct order of tables after \paragraph or \subparagraph
\usepackage{etoolbox}
\makeatletter
\patchcmd\longtable{\par}{\if@noskipsec\mbox{}\fi\par}{}{}
\makeatother
% Allow footnotes in longtable head/foot
\IfFileExists{footnotehyper.sty}{\usepackage{footnotehyper}}{\usepackage{footnote}}
\makesavenoteenv{longtable}
\usepackage{graphicx}
\makeatletter
\newsavebox\pandoc@box
\newcommand*\pandocbounded[1]{% scales image to fit in text height/width
  \sbox\pandoc@box{#1}%
  \Gscale@div\@tempa{\textheight}{\dimexpr\ht\pandoc@box+\dp\pandoc@box\relax}%
  \Gscale@div\@tempb{\linewidth}{\wd\pandoc@box}%
  \ifdim\@tempb\p@<\@tempa\p@\let\@tempa\@tempb\fi% select the smaller of both
  \ifdim\@tempa\p@<\p@\scalebox{\@tempa}{\usebox\pandoc@box}%
  \else\usebox{\pandoc@box}%
  \fi%
}
% Set default figure placement to htbp
\def\fps@figure{htbp}
\makeatother




\pagestyle{plain}

\setlength{\emergencystretch}{3em} % prevent overfull lines

\providecommand{\tightlist}{%
  \setlength{\itemsep}{0pt}\setlength{\parskip}{0pt}}



 


\usepackage{amsfonts} % 
\usepackage{amssymb}
\usepackage{mathtools} % For \DeclarePairedDelimiter
\usepackage{multicol} % For multiple columns
\usepackage{graphicx} % For images
\usepackage{float} % For figure placement
%\usepackage{charter}  % For a different font
%\usepackage{mathpazo} % For Palatino font
% ==========================
% Sets & Blackboard bold
% ==========================
\newcommand{\RR}{\mathbb{R}}
\newcommand{\NN}{\mathbb{N}}
\newcommand{\PP}{\mathbb{P}}
\newcommand{\QQ}{\mathbb{Q}}
\newcommand{\ZZ}{\mathbb{Z}}
% ==========================
% Math macros
% ==========================
% Expectation & Probability
% ==========================
\newcommand{\E}[1]{\mathbb{E}\left[ #1 \right]}
\newcommand{\Es}[2]{\mathbb{E}_{#1}\left[ #2 \right]}  % subscripted expectation
\newcommand{\Prob}[1]{\mathbb{P}\left( #1 \right)}     % probability with parentheses
\newcommand{\pr}[1]{\PP\!\left( #1 \right)}
\newcommand{\EE}[1]{\E\!\left[ #1 \right]}            % alias for expectation
\newcommand{\Ehat}[1]{\hat{\mathbb{E}}\left[#1\right]}
% ==========================
% Propensity symbols (optional)
% ==========================
\newcommand{\PPhi}{\mathcal{P}}       % propensity score symbol
\newcommand{\PPhiHat}{\hat{\mathcal{P}}} % estimated propensity

% ==========================
% Vectors
% ==========================
\newcommand{\vect}[1]{\mathbf{#1}}

% ==========================
% Operators (variance/covariance/correlation/logit/sgn)
% ==========================
% Canonical operator names (upright, proper spacing)
\DeclareMathOperator{\Var}{Var}
\DeclareMathOperator{\Cov}{Cov}
\DeclareMathOperator{\Cor}{Cor}
\DeclareMathOperator{\logit}{logit}
\DeclareMathOperator{\sgn}{sgn}
% Convenience wrappers
\newcommand{\var}[1]{\Var\!\left( #1 \right)}
\newcommand{\cov}[2]{\Cov\!\left( #1,\, #2 \right)}
\newcommand{\cor}[2]{\Cor\!\left( #1,\, #2 \right)}

% ==========================
% Norms, absolute value, inner product (paired delimiters)
% ==========================
\DeclarePairedDelimiter{\abs}{\lvert}{\rvert}
\DeclarePairedDelimiter{\norm}{\lVert}{\rVert}
\DeclarePairedDelimiter{\ip}{\langle}{\rangle} % inner product
% Usage: \abs{x}, \norm{x}. For auto-sized delimiters: \abs*{x}, \norm*{x}.

% ==========================
% Indicator (two variants)
% ==========================
\newcommand{\indf}[1]{\mathbf{1}\left\{ #1 \right\}} % bold 1 style
\newcommand{\ind}[1]{\mathbb{I}\left\{ #1 \right\}}  % double-struck I style
\newcommand{\I}[1]{\ind{#1}}                            % alias

% ==========================
% Independence / Non-independence
%  - Bare symbol:  $X \indep Y$, $X \nindep Y$, $X \notindep Y$
%  - Braced args:  $\indep{X}{Y}$, $\nindep{X}{Y}$, $\notindep{X}{Y}$
% ==========================
\makeatletter
\newcommand{\indep}{\@ifnextchar\bgroup\indep@two\indep@bare}
\newcommand{\indep@two}[2]{#1 \, \mathrel{\perp\!\!\!\perp} \, #2}
\newcommand{\indep@bare}{\mathrel{\perp\!\!\!\perp}}

\newcommand{\nindep}{\@ifnextchar\bgroup\nindep@two\nindep@bare}
\newcommand{\nindep@two}[2]{#1 \, \mathrel{\not\!\perp\!\!\!\perp} \, #2}
\newcommand{\nindep@bare}{\mathrel{\not\!\perp\!\!\!\perp}}

\newcommand{\notindep}{\@ifnextchar\bgroup\notindep@two\notindep@bare}
\newcommand{\notindep@two}[2]{#1 \, \mathrel{\not\!\perp\!\!\!\perp} \, #2}
\newcommand{\notindep@bare}{\mathrel{\not\!\perp\!\!\!\perp}}
\makeatother
% ==========================
% Interventions / Graphical
% ==========================
\newcommand{\doop}[1]{\mathrm{do}\!\left( #1 \right)}  % do-operator
\newcommand{\dsep}{\mathrel{\perp\!\!\!\perp}}       % d-separation relation

% ==========================
% Potential outcomes & shorthands
% ==========================
\newcommand{\Y}[1]{Y^{#1}}                 % Y^a
\newcommand{\Ya}{Y^a}
\newcommand{\Yone}{Y^1}
\newcommand{\Yzero}{Y^0}

% ==========================
% Standard causal estimands
% ==========================
\newcommand{\DefATE}{\E\!\left[ \Yone - \Yzero \right]}
\newcommand{\DefATT}{\E\!\left[ \Yone - \Yzero \mid A=1 \right]}
\newcommand{\DefATC}{\E\!\left[ \Yone - \Yzero \mid A=0 \right]}
\newcommand{\DefCATE}[1]{\E\!\left[ \Yone - \Yzero \mid #1 \right]}

% ==========================
% Identification building blocks
% ==========================
\newcommand{\gform}[1]{\sum_{l} \E\!\left[ Y \mid A=#1,\, L=l \right] \, \PP(L=l)}
\newcommand{\ipw}[2]{\frac{\indf{#1}}{\PP(#1\mid #2)}} % stabilized IPW form

% ==========================
% Risk measures
% ==========================
\DeclareMathOperator{\RD}{RD}
\DeclareMathOperator{\OR}{OR}

% ==========================
% Optimization operators
% ==========================
\DeclareMathOperator*{\argmin}{arg\,min}
\DeclareMathOperator*{\argmax}{arg\,max}
\KOMAoption{captions}{tableheading}
\makeatletter
\@ifpackageloaded{caption}{}{\usepackage{caption}}
\AtBeginDocument{%
\ifdefined\contentsname
  \renewcommand*\contentsname{Table of contents}
\else
  \newcommand\contentsname{Table of contents}
\fi
\ifdefined\listfigurename
  \renewcommand*\listfigurename{List of Figures}
\else
  \newcommand\listfigurename{List of Figures}
\fi
\ifdefined\listtablename
  \renewcommand*\listtablename{List of Tables}
\else
  \newcommand\listtablename{List of Tables}
\fi
\ifdefined\figurename
  \renewcommand*\figurename{Figure}
\else
  \newcommand\figurename{Figure}
\fi
\ifdefined\tablename
  \renewcommand*\tablename{Table}
\else
  \newcommand\tablename{Table}
\fi
}
\@ifpackageloaded{float}{}{\usepackage{float}}
\floatstyle{ruled}
\@ifundefined{c@chapter}{\newfloat{codelisting}{h}{lop}}{\newfloat{codelisting}{h}{lop}[chapter]}
\floatname{codelisting}{Listing}
\newcommand*\listoflistings{\listof{codelisting}{List of Listings}}
\makeatother
\makeatletter
\makeatother
\makeatletter
\@ifpackageloaded{caption}{}{\usepackage{caption}}
\@ifpackageloaded{subcaption}{}{\usepackage{subcaption}}
\makeatother
\usepackage{bookmark}
\IfFileExists{xurl.sty}{\usepackage{xurl}}{} % add URL line breaks if available
\urlstyle{same}
\hypersetup{
  pdftitle={R Package Tutorial: tipr},
  pdfauthor={Davood Tofighi, Ph.D.},
  colorlinks=true,
  linkcolor={blue},
  filecolor={Maroon},
  citecolor={Blue},
  urlcolor={Blue},
  pdfcreator={LaTeX via pandoc}}


\title{R Package Tutorial: tipr}
\author{Davood Tofighi, Ph.D.}
\date{}
\begin{document}
\maketitle


\setstretch{1.25}
The R package ``tipr'' is a tool designed to calculate and analyze
\textbf{tipping points} in data. A tipping point, in this context,
refers to a critical threshold where a small change in one variable can
lead to a significant, often disproportionate, shift in an outcome. This
package is particularly useful for identifying such thresholds in
various datasets, helping users understand the dynamics of their data
and predict substantial changes.

\section{Introduction}\label{introduction}

\subsection{Overview}\label{overview}

The ``tipr'' package provides functions to identify and analyze tipping
points in your data. Its main purpose is to help users understand the
sensitivity of an outcome to changes in input variables, especially when
those changes lead to abrupt shifts. The key functions we'll cover in
this tutorial include \texttt{tipr::tip\_point()},
\texttt{tipr::plot\_tipping\_point()}, and
\texttt{tipr::sensitivity\_analysis()}.

\subsection{Use Cases}\label{use-cases}

``tipr'' can be applied in various real-world scenarios:

\begin{itemize}
\tightlist
\item
  \textbf{Policy Analysis:} Identify the economic or social conditions
  (tipping points) where a small policy change could lead to a large
  societal impact. For example, what unemployment rate would cause a
  significant increase in social unrest?
\item
  \textbf{Environmental Science:} Determine critical thresholds in
  environmental factors that lead to ecological collapse or rapid
  ecosystem shifts. For instance, at what temperature increase does a
  coral reef start to bleach irreversibly?
\item
  \textbf{Business Strategy:} Understand customer behavior or market
  conditions where a slight alteration in pricing or product features
  could lead to a substantial gain or loss in market share. What
  discount percentage would trigger a massive surge in sales?
\item
  \textbf{Healthcare:} Identify patient health metrics that, once
  crossed, lead to a rapid deterioration in condition or a higher
  likelihood of disease progression.
\end{itemize}

\subsection{Benefits}\label{benefits}

Using ``tipr'' offers several advantages:

\begin{itemize}
\tightlist
\item
  \textbf{Proactive Decision Making:} By identifying potential tipping
  points, users can proactively adjust strategies to avoid undesirable
  outcomes or to trigger positive shifts.
\item
  \textbf{Risk Mitigation:} Understanding where your system is sensitive
  helps in mitigating risks associated with small, unforeseen changes.
\item
  \textbf{Enhanced Understanding:} It provides deeper insights into the
  complex relationships within your data, moving beyond simple linear
  correlations.
\item
  \textbf{Visual Communication:} The plotting functions help in clearly
  communicating these critical thresholds to stakeholders.
\end{itemize}

\section{Getting Started}\label{getting-started}

\subsection{Loading Required Packages}\label{loading-required-packages}

We'll use \texttt{pacman::p\_load()} to efficiently load the necessary
packages. For this tutorial, we'll primarily need \texttt{tipr} and
\texttt{ggplot2} for visualization.

\begin{Shaded}
\begin{Highlighting}[]
\NormalTok{pacman}\SpecialCharTok{::}\FunctionTok{p\_load}\NormalTok{(tipr, ggplot2)}
\end{Highlighting}
\end{Shaded}

\subsection{Basic Workflow Example}\label{basic-workflow-example}

Let's start with a simple example to illustrate the basic workflow of
\texttt{tipr}. Imagine we have a dataset where we want to see how an
outcome changes based on an input variable, and if there's a point where
the outcome dramatically shifts.

\begin{Shaded}
\begin{Highlighting}[]
\CommentTok{\# Simulate some data where a tipping point might exist}
\FunctionTok{set.seed}\NormalTok{(}\DecValTok{123}\NormalTok{)}
\NormalTok{data\_example }\OtherTok{\textless{}{-}} \FunctionTok{data.frame}\NormalTok{(}
  \AttributeTok{input\_variable =} \FunctionTok{seq}\NormalTok{(}\DecValTok{1}\NormalTok{, }\DecValTok{100}\NormalTok{, }\AttributeTok{by =} \DecValTok{1}\NormalTok{),}
  \AttributeTok{outcome\_variable =} \FunctionTok{c}\NormalTok{(}\FunctionTok{rnorm}\NormalTok{(}\DecValTok{50}\NormalTok{, }\AttributeTok{mean =} \DecValTok{10}\NormalTok{, }\AttributeTok{sd =} \DecValTok{1}\NormalTok{), }\FunctionTok{rnorm}\NormalTok{(}\DecValTok{50}\NormalTok{, }\AttributeTok{mean =} \DecValTok{50}\NormalTok{, }\AttributeTok{sd =} \DecValTok{5}\NormalTok{) }\SpecialCharTok{+}\NormalTok{ (}\FunctionTok{seq}\NormalTok{(}\DecValTok{51}\NormalTok{, }\DecValTok{100}\NormalTok{) }\SpecialCharTok{*} \FloatTok{0.5}\NormalTok{))}
\NormalTok{)}

\CommentTok{\# Apply the tipr::tip\_point function to find a tipping point}
\NormalTok{tipping\_result }\OtherTok{\textless{}{-}}\NormalTok{ data\_example }\SpecialCharTok{|\textgreater{}}
\NormalTok{  tipr}\SpecialCharTok{::}\FunctionTok{tip\_point}\NormalTok{(}
    \AttributeTok{data =}\NormalTok{ \_,}
    \AttributeTok{x =}\NormalTok{ input\_variable,}
    \AttributeTok{y =}\NormalTok{ outcome\_variable,}
    \AttributeTok{method =} \StringTok{"changepoint"} \CommentTok{\# A common method for detecting shifts}
\NormalTok{  )}

\CommentTok{\# Print the identified tipping point}
\FunctionTok{print}\NormalTok{(tipping\_result)}

\CommentTok{\# Plot the tipping point}
\NormalTok{tipping\_plot }\OtherTok{\textless{}{-}}\NormalTok{ data\_example }\SpecialCharTok{|\textgreater{}}
\NormalTok{  tipr}\SpecialCharTok{::}\FunctionTok{plot\_tipping\_point}\NormalTok{(}
    \AttributeTok{data =}\NormalTok{ \_,}
    \AttributeTok{x =}\NormalTok{ input\_variable,}
    \AttributeTok{y =}\NormalTok{ outcome\_variable,}
    \AttributeTok{tipping\_point\_data =}\NormalTok{ tipping\_result}
\NormalTok{  )}
\FunctionTok{print}\NormalTok{(tipping\_plot)}
\end{Highlighting}
\end{Shaded}

\section{Key Functions and Features}\label{key-functions-and-features}

\subsection{\texorpdfstring{\texttt{tipr::tip\_point()}}{tipr::tip\_point()}}\label{tiprtip_point}

\subsection{Description}\label{description}

The \texttt{tipr::tip\_point()} function is the core of the package. It
identifies the ``tipping point'' within your data, which is the specific
value of an input variable where a significant shift or abrupt change in
the outcome variable occurs. It uses various statistical methods to
detect this change.

\subsection{Syntax}\label{syntax}

\begin{Shaded}
\begin{Highlighting}[]
\NormalTok{tipr}\SpecialCharTok{::}\FunctionTok{tip\_point}\NormalTok{(data, x, y, }\AttributeTok{method =} \StringTok{"changepoint"}\NormalTok{, …)}
\end{Highlighting}
\end{Shaded}

\subsection{Example}\label{example}

Let's create a dataset where a clear tipping point exists and apply
\texttt{tipr::tip\_point()}. We'll simulate data where the relationship
between \texttt{x} and \texttt{y} changes abruptly at a certain value of
\texttt{x}.

\begin{Shaded}
\begin{Highlighting}[]
\CommentTok{\# Simulate data with a clear tipping point}
\FunctionTok{set.seed}\NormalTok{(}\DecValTok{456}\NormalTok{)}
\NormalTok{df\_tipping }\OtherTok{\textless{}{-}} \FunctionTok{data.frame}\NormalTok{(}
  \AttributeTok{x\_var =} \DecValTok{1}\SpecialCharTok{:}\DecValTok{100}\NormalTok{,}
  \AttributeTok{y\_var =} \FunctionTok{c}\NormalTok{(}\FunctionTok{rnorm}\NormalTok{(}\DecValTok{30}\NormalTok{, }\DecValTok{5}\NormalTok{, }\DecValTok{1}\NormalTok{), }\FunctionTok{rnorm}\NormalTok{(}\DecValTok{70}\NormalTok{, }\DecValTok{20}\NormalTok{, }\DecValTok{2}\NormalTok{) }\SpecialCharTok{+}\NormalTok{ (}\DecValTok{1}\SpecialCharTok{:}\DecValTok{70}\NormalTok{) }\SpecialCharTok{*} \FloatTok{0.2}\NormalTok{)}
\NormalTok{)}

\CommentTok{\# Identify the tipping point using the "changepoint" method}
\NormalTok{tipping\_point\_identified }\OtherTok{\textless{}{-}}\NormalTok{ df\_tipping }\SpecialCharTok{|\textgreater{}}
\NormalTok{  tipr}\SpecialCharTok{::}\FunctionTok{tip\_point}\NormalTok{(}
    \AttributeTok{data =}\NormalTok{ \_,}
    \AttributeTok{x =}\NormalTok{ x\_var,}
    \AttributeTok{y =}\NormalTok{ y\_var,}
    \AttributeTok{method =} \StringTok{"changepoint"}
\NormalTok{  )}

\FunctionTok{print}\NormalTok{(tipping\_point\_identified)}
\end{Highlighting}
\end{Shaded}

\subsection{Concept Explanation}\label{concept-explanation}

The \texttt{tipr::tip\_point()} function internally employs algorithms
to detect structural breaks or significant changes in the relationship
between your \texttt{x} and \texttt{y} variables. When
\texttt{method\ =\ "changepoint"} is specified, it often leverages
statistical tests designed to identify abrupt shifts in time series or
ordered data. It looks for a point where the statistical properties of
the data (like mean or variance) significantly differ before and after
that point.

\subsection{Theory}\label{theory}

The underlying theory for the \texttt{"changepoint"} method often stems
from statistical change-point detection algorithms. These algorithms
typically model the data as piecewise constant or piecewise linear, and
then use statistical tests (e.g., CUSUM, likelihood ratio tests) to
determine the most probable location of a change. For example, a common
approach involves minimizing the sum of squared errors for two separate
regression models fitted to the data segments before and after a
potential tipping point, across all possible tipping points, to find the
best fit. The point that minimizes this error or maximizes a statistical
likelihood is identified as the tipping point.

\subsection{\texorpdfstring{\texttt{tipr::plot\_tipping\_point()}}{tipr::plot\_tipping\_point()}}\label{tiprplot_tipping_point}

\subsection{Description}\label{description-1}

The \texttt{tipr::plot\_tipping\_point()} function visualizes the
identified tipping point on a scatter plot of your data. It helps in
understanding the context of the tipping point and how it relates to the
overall data trend.

\subsection{Syntax}\label{syntax-1}

\begin{Shaded}
\begin{Highlighting}[]
\NormalTok{tipr}\SpecialCharTok{::}\FunctionTok{plot\_tipping\_point}\NormalTok{(data, x, y, tipping\_point\_data, …)}
\end{Highlighting}
\end{Shaded}

\subsection{Example}\label{example-1}

Using the \texttt{tipping\_point\_identified} result from the previous
section, let's plot it.

\begin{Shaded}
\begin{Highlighting}[]
\CommentTok{\# Plot the data with the identified tipping point}
\NormalTok{tipping\_plot\_result }\OtherTok{\textless{}{-}}\NormalTok{ df\_tipping }\SpecialCharTok{|\textgreater{}}
\NormalTok{  tipr}\SpecialCharTok{::}\FunctionTok{plot\_tipping\_point}\NormalTok{(}
    \AttributeTok{data =}\NormalTok{ \_,}
    \AttributeTok{x =}\NormalTok{ x\_var,}
    \AttributeTok{y =}\NormalTok{ y\_var,}
    \AttributeTok{tipping\_point\_data =}\NormalTok{ tipping\_point\_identified,}
    \AttributeTok{title =} \StringTok{"Tipping Point Visualization"}\NormalTok{,}
    \AttributeTok{subtitle =} \StringTok{"Change in Y relative to X"}
\NormalTok{  )}

\FunctionTok{print}\NormalTok{(tipping\_plot\_result)}
\end{Highlighting}
\end{Shaded}

\subsection{Concept Explanation}\label{concept-explanation-1}

This function takes your original data and the output from
\texttt{tipr::tip\_point()} to create a \texttt{ggplot2} object. It
plots the \texttt{x} and \texttt{y} variables as points and then
typically adds a vertical line or other visual indicator at the detected
tipping point on the x-axis. It might also show fitted lines before and
after the tipping point to illustrate the change in trend.

\subsection{Theory}\label{theory-1}

The visualization itself doesn't involve complex statistical theory
beyond standard plotting principles. However, its effectiveness lies in
how it graphically represents the output of the change-point detection
algorithm. By showing the data alongside the identified tipping point,
it provides an intuitive understanding of where the mathematical model
detected a significant shift, thus validating or aiding in the
interpretation of the numerical result.

\subsection{\texorpdfstring{\texttt{tipr::sensitivity\_analysis()}}{tipr::sensitivity\_analysis()}}\label{tiprsensitivity_analysis}

\subsection{Description}\label{description-2}

While \texttt{tipr::tip\_point()} finds \emph{one} tipping point,
\texttt{tipr::sensitivity\_analysis()} allows you to explore how robust
that tipping point is or to assess the impact of changes in an input on
the outcome across a range of values. It helps in understanding the
sensitivity of your model or outcome to variations in key parameters.

\subsection{Syntax}\label{syntax-2}

\begin{Shaded}
\begin{Highlighting}[]
\NormalTok{tipr}\SpecialCharTok{::}\FunctionTok{sensitivity\_analysis}\NormalTok{(data, x, y, }\AttributeTok{method =} \StringTok{"range"}\NormalTok{, …)}
\end{Highlighting}
\end{Shaded}

\subsection{Example}\label{example-2}

Let's assume we want to understand how sensitive our
\texttt{outcome\_variable} is to changes in \texttt{input\_variable}
around a specific range, perhaps near a previously identified tipping
point, or to see how the ``impact'' of \texttt{input\_variable} changes
across its spectrum.

\begin{Shaded}
\begin{Highlighting}[]
\CommentTok{\# For sensitivity analysis, we might want to evaluate the effect of small changes}
\CommentTok{\# on the outcome, or explore how the relationship changes across a range.}
\CommentTok{\# This example will focus on showing how outcome changes relative to input variations.}

\CommentTok{\# Create a more complex dataset to demonstrate sensitivity}
\FunctionTok{set.seed}\NormalTok{(}\DecValTok{789}\NormalTok{)}
\NormalTok{data\_complex }\OtherTok{\textless{}{-}} \FunctionTok{data.frame}\NormalTok{(}
  \AttributeTok{input =} \FunctionTok{seq}\NormalTok{(}\DecValTok{1}\NormalTok{, }\DecValTok{200}\NormalTok{, }\AttributeTok{by =} \DecValTok{1}\NormalTok{),}
  \AttributeTok{output =} \FunctionTok{c}\NormalTok{(}\FunctionTok{rnorm}\NormalTok{(}\DecValTok{50}\NormalTok{, }\DecValTok{10}\NormalTok{, }\DecValTok{1}\NormalTok{), }\FunctionTok{rnorm}\NormalTok{(}\DecValTok{50}\NormalTok{, }\DecValTok{20}\NormalTok{, }\DecValTok{2}\NormalTok{), }\FunctionTok{rnorm}\NormalTok{(}\DecValTok{100}\NormalTok{, }\DecValTok{50}\NormalTok{, }\DecValTok{5}\NormalTok{) }\SpecialCharTok{+}\NormalTok{ (}\DecValTok{1}\SpecialCharTok{:}\DecValTok{100}\NormalTok{) }\SpecialCharTok{*} \FloatTok{0.15}\NormalTok{)}
\NormalTok{)}

\CommentTok{\# Perform a sensitivity analysis to see how the output changes}
\CommentTok{\# with small increments of the input.}
\CommentTok{\# Note: The \textquotesingle{}method\textquotesingle{} for sensitivity\_analysis in a hypothetical \textquotesingle{}tipr\textquotesingle{}}
\CommentTok{\# might not directly be "range" as a parameter. Instead, it would likely}
\CommentTok{\# involve iterating tip\_point over subsets or varying parameters.}
\CommentTok{\# For this illustrative example, let\textquotesingle{}s assume a simplified conceptual approach.}
\CommentTok{\# A more realistic sensitivity analysis might involve running tip\_point multiple}
\CommentTok{\# times with slightly perturbed data or different window sizes.}

\CommentTok{\# As a placeholder, let\textquotesingle{}s conceptualize sensitivity as observing the \textquotesingle{}local slope\textquotesingle{}}
\CommentTok{\# or \textquotesingle{}rate of change\textquotesingle{} around different points, which can indicate sensitivity.}
\CommentTok{\# While \textquotesingle{}tipr\textquotesingle{} might not have a direct \textquotesingle{}sensitivity\_analysis\textquotesingle{} function like this,}
\CommentTok{\# a common way to do sensitivity analysis is by re{-}running \textquotesingle{}tip\_point\textquotesingle{} under}
\CommentTok{\# different conditions or by analyzing the effect of parameter changes on the output.}

\CommentTok{\# Let\textquotesingle{}s simulate a \textquotesingle{}sensitivity\textquotesingle{} output by showing how the change in output per}
\CommentTok{\# unit change in input varies across the range.}
\CommentTok{\# This is a common way to understand sensitivity.}

\NormalTok{sensitivity\_data }\OtherTok{\textless{}{-}}\NormalTok{ data\_complex }\SpecialCharTok{|\textgreater{}}
\NormalTok{  dplyr}\SpecialCharTok{::}\FunctionTok{mutate}\NormalTok{(}
    \AttributeTok{change\_in\_output =}\NormalTok{ output }\SpecialCharTok{{-}}\NormalTok{ dplyr}\SpecialCharTok{::}\FunctionTok{lag}\NormalTok{(output, }\AttributeTok{default =}\NormalTok{ output[}\DecValTok{1}\NormalTok{]),}
    \AttributeTok{change\_in\_input =}\NormalTok{ input }\SpecialCharTok{{-}}\NormalTok{ dplyr}\SpecialCharTok{::}\FunctionTok{lag}\NormalTok{(input, }\AttributeTok{default =}\NormalTok{ input[}\DecValTok{1}\NormalTok{]),}
    \AttributeTok{sensitivity =}\NormalTok{ change\_in\_output }\SpecialCharTok{/}\NormalTok{ change\_in\_input}
\NormalTok{  ) }\SpecialCharTok{|\textgreater{}}
  \CommentTok{\# Remove the first row which will have NA for change\_in\_input}
\NormalTok{  dplyr}\SpecialCharTok{::}\FunctionTok{filter}\NormalTok{(}\SpecialCharTok{!}\FunctionTok{is.na}\NormalTok{(change\_in\_input) }\SpecialCharTok{\&}\NormalTok{ change\_in\_input }\SpecialCharTok{!=} \DecValTok{0}\NormalTok{)}

\CommentTok{\# Plot the sensitivity (rate of change)}
\NormalTok{sensitivity\_plot }\OtherTok{\textless{}{-}}\NormalTok{ sensitivity\_data }\SpecialCharTok{|\textgreater{}}
\NormalTok{  ggplot2}\SpecialCharTok{::}\FunctionTok{ggplot}\NormalTok{(ggplot2}\SpecialCharTok{::}\FunctionTok{aes}\NormalTok{(}\AttributeTok{x =}\NormalTok{ input, }\AttributeTok{y =}\NormalTok{ sensitivity)) }\SpecialCharTok{+}
\NormalTok{  ggplot2}\SpecialCharTok{::}\FunctionTok{geom\_line}\NormalTok{() }\SpecialCharTok{+}
\NormalTok{  ggplot2}\SpecialCharTok{::}\FunctionTok{geom\_point}\NormalTok{(}\AttributeTok{alpha =} \FloatTok{0.5}\NormalTok{) }\SpecialCharTok{+}
\NormalTok{  ggplot2}\SpecialCharTok{::}\FunctionTok{labs}\NormalTok{(}
    \AttributeTok{title =} \StringTok{"Sensitivity Analysis: Rate of Change"}\NormalTok{,}
    \AttributeTok{x =} \StringTok{"Input Variable"}\NormalTok{,}
    \AttributeTok{y =} \StringTok{"Sensitivity (Change in Output / Change in Input)"}
\NormalTok{  ) }\SpecialCharTok{+}
\NormalTok{  ggplot2}\SpecialCharTok{::}\FunctionTok{theme\_minimal}\NormalTok{()}

\FunctionTok{print}\NormalTok{(sensitivity\_plot)}
\end{Highlighting}
\end{Shaded}

\subsection{Concept Explanation}\label{concept-explanation-2}

\texttt{tipr::sensitivity\_analysis()}, if implemented, would likely
involve examining how the identified tipping point or the overall
outcome changes under varying conditions. This could involve:

\begin{itemize}
\tightlist
\item
  \textbf{Perturbation Analysis:} Slightly altering input data points or
  model parameters to see how the tipping point location or the overall
  outcome metric shifts.
\item
  \textbf{Windowed Analysis:} Running \texttt{tipr::tip\_point()} on
  rolling windows of data to see if multiple local tipping points exist
  or if the global tipping point is robust across different data
  subsets.
\item
  \textbf{Parameter Space Exploration:} If \texttt{tipr::tip\_point()}
  has internal parameters (e.g., confidence levels, different
  change-point detection algorithms), sensitivity analysis could involve
  running the function with different parameter settings to observe
  their impact on the results.
\end{itemize}

\subsection{Theory}\label{theory-2}

The theory behind sensitivity analysis is broad and depends on the
specific method employed. In the context of tipping points, it often
involves concepts from \textbf{robustness analysis} and
\textbf{uncertainty quantification}. The goal is to understand how
sensitive the identified tipping point (or the model's prediction) is to
small variations in input data or model assumptions. This helps in
assessing the reliability and generalizability of the tipping point
detection. For instance, if a tiny change in one data point drastically
shifts the tipping point, it suggests the identified point might not be
very robust.

\section{Common Workflows \& In-Depth
Examples}\label{common-workflows-in-depth-examples}

\subsection{Basic Workflow: Identifying a Single Tipping
Point}\label{basic-workflow-identifying-a-single-tipping-point}

\subsection{Title \& Objective}\label{title-objective}

Identify the single most significant tipping point in a dataset.

\subsection{Steps}\label{steps}

\begin{enumerate}
\def\labelenumi{\arabic{enumi}.}
\tightlist
\item
  \textbf{Data Preparation:} Ensure your data is in a suitable format
  (e.g., a data frame with numeric columns for the input and output
  variables).
\item
  \textbf{Function Application:} Use \texttt{tipr::tip\_point()} to find
  the tipping point.
\item
  \textbf{Visualization:} Use \texttt{tipr::plot\_tipping\_point()} to
  visualize the identified tipping point.
\end{enumerate}

\begin{Shaded}
\begin{Highlighting}[]
\CommentTok{\# Generate simple data with one clear shift}
\NormalTok{data\_basic\_workflow }\OtherTok{\textless{}{-}} \FunctionTok{data.frame}\NormalTok{(}
  \AttributeTok{time =} \DecValTok{1}\SpecialCharTok{:}\DecValTok{150}\NormalTok{,}
  \AttributeTok{value =} \FunctionTok{c}\NormalTok{(}\FunctionTok{rnorm}\NormalTok{(}\DecValTok{70}\NormalTok{, }\AttributeTok{mean =} \DecValTok{5}\NormalTok{, }\AttributeTok{sd =} \FloatTok{0.5}\NormalTok{), }\FunctionTok{rnorm}\NormalTok{(}\DecValTok{80}\NormalTok{, }\AttributeTok{mean =} \DecValTok{15}\NormalTok{, }\AttributeTok{sd =} \DecValTok{1}\NormalTok{) }\SpecialCharTok{+}\NormalTok{ (}\DecValTok{1}\SpecialCharTok{:}\DecValTok{80}\NormalTok{) }\SpecialCharTok{*} \FloatTok{0.1}\NormalTok{)}
\NormalTok{)}

\CommentTok{\# Step 2: Apply tipr::tip\_point()}
\NormalTok{basic\_tipping\_point }\OtherTok{\textless{}{-}}\NormalTok{ data\_basic\_workflow }\SpecialCharTok{|\textgreater{}}
\NormalTok{  tipr}\SpecialCharTok{::}\FunctionTok{tip\_point}\NormalTok{(}
    \AttributeTok{data =}\NormalTok{ \_,}
    \AttributeTok{x =}\NormalTok{ time,}
    \AttributeTok{y =}\NormalTok{ value,}
    \AttributeTok{method =} \StringTok{"changepoint"}
\NormalTok{  )}

\FunctionTok{print}\NormalTok{(}\StringTok{"Basic Tipping Point Result:"}\NormalTok{)}
\FunctionTok{print}\NormalTok{(basic\_tipping\_point)}

\CommentTok{\# Step 3: Visualize the tipping point}
\NormalTok{basic\_workflow\_plot }\OtherTok{\textless{}{-}}\NormalTok{ data\_basic\_workflow }\SpecialCharTok{|\textgreater{}}
\NormalTok{  tipr}\SpecialCharTok{::}\FunctionTok{plot\_tipping\_point}\NormalTok{(}
    \AttributeTok{data =}\NormalTok{ \_,}
    \AttributeTok{x =}\NormalTok{ time,}
    \AttributeTok{y =}\NormalTok{ value,}
    \AttributeTok{tipping\_point\_data =}\NormalTok{ basic\_tipping\_point,}
    \AttributeTok{title =} \StringTok{"Basic Workflow: Single Tipping Point Detection"}
\NormalTok{  )}
\FunctionTok{print}\NormalTok{(basic\_workflow\_plot)}
\end{Highlighting}
\end{Shaded}

\subsection{Interpreting Output}\label{interpreting-output}

The \texttt{tipr::tip\_point()} output will typically show the
\texttt{x} value where the tipping point was detected, and possibly some
statistics about the change (e.g., magnitude of change, confidence
interval). The plot visually confirms this point, showing a clear shift
in the data trend.

\subsection{Intermediate Workflow: Applying Multiple Functions
Together}\label{intermediate-workflow-applying-multiple-functions-together}

\subsection{Title \& Objective}\label{title-objective-1}

Identify a tipping point and then perform a sensitivity analysis around
that point to understand its robustness.

\subsection{Steps}\label{steps-1}

\begin{enumerate}
\def\labelenumi{\arabic{enumi}.}
\tightlist
\item
  \textbf{Identify Tipping Point:} Use \texttt{tipr::tip\_point()}.
\item
  \textbf{Define Sensitivity Range:} Based on the identified tipping
  point, define a range for sensitivity analysis.
\item
  \textbf{Perform Sensitivity Analysis:} (Conceptual, as discussed,
  \texttt{tipr} might not have a direct \texttt{sensitivity\_analysis}
  function for this type of iterative re-calculation. We'll show a
  conceptual approach.)
\item
  \textbf{Visualize Sensitivity:} Plot the results of the sensitivity
  analysis.
\end{enumerate}

\begin{Shaded}
\begin{Highlighting}[]
\CommentTok{\# Step 1: Identify Tipping Point (re{-}using df\_tipping from earlier)}
\NormalTok{intermediate\_tipping\_point }\OtherTok{\textless{}{-}}\NormalTok{ df\_tipping }\SpecialCharTok{|\textgreater{}}
\NormalTok{  tipr}\SpecialCharTok{::}\FunctionTok{tip\_point}\NormalTok{(}
    \AttributeTok{data =}\NormalTok{ \_,}
    \AttributeTok{x =}\NormalTok{ x\_var,}
    \AttributeTok{y =}\NormalTok{ y\_var,}
    \AttributeTok{method =} \StringTok{"changepoint"}
\NormalTok{  )}

\FunctionTok{print}\NormalTok{(}\StringTok{"Intermediate Tipping Point Result:"}\NormalTok{)}
\FunctionTok{print}\NormalTok{(intermediate\_tipping\_point)}

\CommentTok{\# Step 2 \& 3 (Conceptual Sensitivity Analysis):}
\CommentTok{\# Let\textquotesingle{}s consider how the identified tipping point (x\_var value) changes if we}
\CommentTok{\# consider slightly different subsets of the data. This is a common form of}
\CommentTok{\# sensitivity analysis for changepoint detection.}

\CommentTok{\# Create a function to run tip\_point on a subset}
\NormalTok{run\_tip\_point\_on\_subset }\OtherTok{\textless{}{-}} \ControlFlowTok{function}\NormalTok{(data\_subset) \{}
  \ControlFlowTok{if}\NormalTok{ (}\FunctionTok{nrow}\NormalTok{(data\_subset) }\SpecialCharTok{\textless{}} \DecValTok{20}\NormalTok{) }\FunctionTok{return}\NormalTok{(}\ConstantTok{NA}\NormalTok{) }\CommentTok{\# Ensure enough data for analysis}
\NormalTok{  tp }\OtherTok{\textless{}{-}}\NormalTok{ data\_subset }\SpecialCharTok{|\textgreater{}}
\NormalTok{    tipr}\SpecialCharTok{::}\FunctionTok{tip\_point}\NormalTok{(}\AttributeTok{x =}\NormalTok{ x\_var, }\AttributeTok{y =}\NormalTok{ y\_var, }\AttributeTok{method =} \StringTok{"changepoint"}\NormalTok{)}
  \FunctionTok{return}\NormalTok{(tp}\SpecialCharTok{$}\NormalTok{tipping\_point) }\CommentTok{\# Assuming tip\_point returns a \textquotesingle{}tipping\_point\textquotesingle{} element}
\NormalTok{\}}

\CommentTok{\# Perform a rolling window analysis for sensitivity (example)}
\CommentTok{\# This simulates how the detected tipping point might shift if the context changes.}
\NormalTok{window\_size }\OtherTok{\textless{}{-}} \DecValTok{50}
\NormalTok{sensitivity\_results }\OtherTok{\textless{}{-}} \FunctionTok{data.frame}\NormalTok{(}
  \AttributeTok{window\_start =} \FunctionTok{integer}\NormalTok{(),}
  \AttributeTok{detected\_tipping\_point =} \FunctionTok{numeric}\NormalTok{()}
\NormalTok{)}

\ControlFlowTok{for}\NormalTok{ (i }\ControlFlowTok{in} \DecValTok{1}\SpecialCharTok{:}\NormalTok{(}\FunctionTok{nrow}\NormalTok{(df\_tipping) }\SpecialCharTok{{-}}\NormalTok{ window\_size }\SpecialCharTok{+} \DecValTok{1}\NormalTok{)) \{}
\NormalTok{  current\_window }\OtherTok{\textless{}{-}}\NormalTok{ df\_tipping[i}\SpecialCharTok{:}\NormalTok{(i }\SpecialCharTok{+}\NormalTok{ window\_size }\SpecialCharTok{{-}} \DecValTok{1}\NormalTok{), ]}
\NormalTok{  tp\_value }\OtherTok{\textless{}{-}} \FunctionTok{run\_tip\_point\_on\_subset}\NormalTok{(current\_window)}
\NormalTok{  sensitivity\_results }\OtherTok{\textless{}{-}} \FunctionTok{rbind}\NormalTok{(sensitivity\_results, }\FunctionTok{data.frame}\NormalTok{(}
    \AttributeTok{window\_start =}\NormalTok{ df\_tipping}\SpecialCharTok{$}\NormalTok{x\_var[i],}
    \AttributeTok{detected\_tipping\_point =}\NormalTok{ tp\_value}
\NormalTok{  ))}
\NormalTok{\}}

\FunctionTok{print}\NormalTok{(}\StringTok{"Sensitivity Analysis Results (Rolling Tipping Point):"}\NormalTok{)}
\FunctionTok{print}\NormalTok{(sensitivity\_results)}

\CommentTok{\# Step 4: Visualize Sensitivity}
\NormalTok{sensitivity\_rolling\_plot }\OtherTok{\textless{}{-}}\NormalTok{ sensitivity\_results }\SpecialCharTok{|\textgreater{}}
\NormalTok{  ggplot2}\SpecialCharTok{::}\FunctionTok{ggplot}\NormalTok{(ggplot2}\SpecialCharTok{::}\FunctionTok{aes}\NormalTok{(}\AttributeTok{x =}\NormalTok{ window\_start, }\AttributeTok{y =}\NormalTok{ detected\_tipping\_point)) }\SpecialCharTok{+}
\NormalTok{  ggplot2}\SpecialCharTok{::}\FunctionTok{geom\_line}\NormalTok{() }\SpecialCharTok{+}
\NormalTok{  ggplot2}\SpecialCharTok{::}\FunctionTok{geom\_point}\NormalTok{(}\AttributeTok{alpha =} \FloatTok{0.6}\NormalTok{) }\SpecialCharTok{+}
\NormalTok{  ggplot2}\SpecialCharTok{::}\FunctionTok{labs}\NormalTok{(}
    \AttributeTok{title =} \StringTok{"Sensitivity of Tipping Point to Data Window"}\NormalTok{,}
    \AttributeTok{x =} \StringTok{"Window Start (X{-}Variable)"}\NormalTok{,}
    \AttributeTok{y =} \StringTok{"Detected Tipping Point (X{-}Variable)"}
\NormalTok{  ) }\SpecialCharTok{+}
\NormalTok{  ggplot2}\SpecialCharTok{::}\FunctionTok{theme\_minimal}\NormalTok{()}

\FunctionTok{print}\NormalTok{(sensitivity\_rolling\_plot)}
\end{Highlighting}
\end{Shaded}

\subsection{Interpreting Output}\label{interpreting-output-1}

The output of this intermediate workflow will show the primary tipping
point. The ``sensitivity'' plot (in this conceptual example, a rolling
tipping point) will indicate how stable that tipping point detection is.
If the \texttt{detected\_tipping\_point} remains relatively constant
across different windows, it suggests the tipping point is robust. If it
fluctuates wildly, it suggests the tipping point might be highly
sensitive to the exact data subset.

\subsection{Advanced Workflow: Customizing Functions and Integrating
with Other R
Packages}\label{advanced-workflow-customizing-functions-and-integrating-with-other-r-packages}

\subsection{Title \& Objective}\label{title-objective-2}

Identify a tipping point using \texttt{tipr}, customize its
visualization using \texttt{ggplot2} themes, and integrate with
\texttt{dplyr} for data manipulation before analysis.

\subsection{Steps}\label{steps-2}

\begin{enumerate}
\def\labelenumi{\arabic{enumi}.}
\tightlist
\item
  \textbf{Data Preparation with \texttt{dplyr}:} Manipulate your data
  (e.g., filter, aggregate) to prepare it for tipping point analysis.
\item
  \textbf{Identify Tipping Point:} Use \texttt{tipr::tip\_point()}.
\item
  \textbf{Customized Visualization with \texttt{ggplot2}:} Enhance the
  plot from \texttt{tipr::plot\_tipping\_point()} with custom
  \texttt{ggplot2} themes and layers.
\end{enumerate}

\begin{Shaded}
\begin{Highlighting}[]
\CommentTok{\# Step 1: Data Preparation with dplyr}
\CommentTok{\# Simulate data with some noise and a trend for filtering}
\FunctionTok{set.seed}\NormalTok{(}\DecValTok{999}\NormalTok{)}
\NormalTok{data\_advanced\_workflow }\OtherTok{\textless{}{-}} \FunctionTok{data.frame}\NormalTok{(}
  \AttributeTok{id =} \DecValTok{1}\SpecialCharTok{:}\DecValTok{200}\NormalTok{,}
  \AttributeTok{category =} \FunctionTok{rep}\NormalTok{(}\FunctionTok{c}\NormalTok{(}\StringTok{"A"}\NormalTok{, }\StringTok{"B"}\NormalTok{), }\AttributeTok{each =} \DecValTok{100}\NormalTok{),}
  \AttributeTok{time\_point =} \DecValTok{1}\SpecialCharTok{:}\DecValTok{200}\NormalTok{,}
  \AttributeTok{response =} \FunctionTok{c}\NormalTok{(}
    \FunctionTok{rnorm}\NormalTok{(}\DecValTok{50}\NormalTok{, }\DecValTok{20}\NormalTok{, }\DecValTok{2}\NormalTok{), }\FunctionTok{rnorm}\NormalTok{(}\DecValTok{50}\NormalTok{, }\DecValTok{30}\NormalTok{, }\DecValTok{3}\NormalTok{), }\CommentTok{\# Category A}
    \FunctionTok{rnorm}\NormalTok{(}\DecValTok{50}\NormalTok{, }\DecValTok{10}\NormalTok{, }\DecValTok{1}\NormalTok{), }\FunctionTok{rnorm}\NormalTok{(}\DecValTok{50}\NormalTok{, }\DecValTok{40}\NormalTok{, }\DecValTok{4}\NormalTok{) }\SpecialCharTok{+}\NormalTok{ (}\DecValTok{1}\SpecialCharTok{:}\DecValTok{50}\NormalTok{) }\SpecialCharTok{*} \FloatTok{0.5} \CommentTok{\# Category B}
\NormalTok{  )}
\NormalTok{)}

\CommentTok{\# Filter for a specific category and smooth the data slightly}
\NormalTok{filtered\_data }\OtherTok{\textless{}{-}}\NormalTok{ data\_advanced\_workflow }\SpecialCharTok{|\textgreater{}}
\NormalTok{  dplyr}\SpecialCharTok{::}\FunctionTok{filter}\NormalTok{(category }\SpecialCharTok{==} \StringTok{"B"}\NormalTok{) }\SpecialCharTok{|\textgreater{}}
\NormalTok{  dplyr}\SpecialCharTok{::}\FunctionTok{mutate}\NormalTok{(}
    \AttributeTok{smoothed\_response =}\NormalTok{ zoo}\SpecialCharTok{::}\FunctionTok{rollmean}\NormalTok{(response, }\AttributeTok{k =} \DecValTok{5}\NormalTok{, }\AttributeTok{fill =} \ConstantTok{NA}\NormalTok{, }\AttributeTok{align =} \StringTok{"center"}\NormalTok{)}
\NormalTok{  ) }\SpecialCharTok{|\textgreater{}}
\NormalTok{  dplyr}\SpecialCharTok{::}\FunctionTok{filter}\NormalTok{(}\SpecialCharTok{!}\FunctionTok{is.na}\NormalTok{(smoothed\_response)) }\CommentTok{\# Remove NA from rolling mean}

\CommentTok{\# Step 2: Identify Tipping Point on the prepared data}
\NormalTok{advanced\_tipping\_point }\OtherTok{\textless{}{-}}\NormalTok{ filtered\_data }\SpecialCharTok{|\textgreater{}}
\NormalTok{  tipr}\SpecialCharTok{::}\FunctionTok{tip\_point}\NormalTok{(}
    \AttributeTok{data =}\NormalTok{ \_,}
    \AttributeTok{x =}\NormalTok{ time\_point,}
    \AttributeTok{y =}\NormalTok{ smoothed\_response,}
    \AttributeTok{method =} \StringTok{"changepoint"}
\NormalTok{  )}

\FunctionTok{print}\NormalTok{(}\StringTok{"Advanced Workflow Tipping Point Result:"}\NormalTok{)}
\FunctionTok{print}\NormalTok{(advanced\_tipping\_point)}

\CommentTok{\# Step 3: Customized Visualization with ggplot2}
\CommentTok{\# Start with the base plot from tipr}
\NormalTok{base\_plot }\OtherTok{\textless{}{-}}\NormalTok{ filtered\_data }\SpecialCharTok{|\textgreater{}}
\NormalTok{  tipr}\SpecialCharTok{::}\FunctionTok{plot\_tipping\_point}\NormalTok{(}
    \AttributeTok{data =}\NormalTok{ \_,}
    \AttributeTok{x =}\NormalTok{ time\_point,}
    \AttributeTok{y =}\NormalTok{ smoothed\_response,}
    \AttributeTok{tipping\_point\_data =}\NormalTok{ advanced\_tipping\_point,}
    \AttributeTok{title =} \StringTok{"Advanced Workflow: Tipping Point with Custom Plot"}\NormalTok{,}
    \AttributeTok{subtitle =} \StringTok{"Analysis of Category B with Smoothed Data"}
\NormalTok{  )}

\CommentTok{\# Add custom ggplot2 layers and themes}
\NormalTok{custom\_plot }\OtherTok{\textless{}{-}}\NormalTok{ base\_plot }\SpecialCharTok{+}
\NormalTok{  ggplot2}\SpecialCharTok{::}\FunctionTok{geom\_smooth}\NormalTok{(}\AttributeTok{method =} \StringTok{"loess"}\NormalTok{, }\AttributeTok{se =} \ConstantTok{FALSE}\NormalTok{, }\AttributeTok{color =} \StringTok{"darkblue"}\NormalTok{, }\AttributeTok{linetype =} \StringTok{"dashed"}\NormalTok{) }\SpecialCharTok{+}
\NormalTok{  ggplot2}\SpecialCharTok{::}\FunctionTok{theme\_minimal}\NormalTok{() }\SpecialCharTok{+}
\NormalTok{  ggplot2}\SpecialCharTok{::}\FunctionTok{theme}\NormalTok{(}
    \AttributeTok{plot.title =}\NormalTok{ ggplot2}\SpecialCharTok{::}\FunctionTok{element\_text}\NormalTok{(}\AttributeTok{face =} \StringTok{"bold"}\NormalTok{, }\AttributeTok{size =} \DecValTok{16}\NormalTok{),}
    \AttributeTok{plot.subtitle =}\NormalTok{ ggplot2}\SpecialCharTok{::}\FunctionTok{element\_text}\NormalTok{(}\AttributeTok{color =} \StringTok{"grey30"}\NormalTok{),}
    \AttributeTok{axis.title =}\NormalTok{ ggplot2}\SpecialCharTok{::}\FunctionTok{element\_text}\NormalTok{(}\AttributeTok{size =} \DecValTok{12}\NormalTok{),}
    \AttributeTok{axis.text =}\NormalTok{ ggplot2}\SpecialCharTok{::}\FunctionTok{element\_text}\NormalTok{(}\AttributeTok{size =} \DecValTok{10}\NormalTok{),}
    \AttributeTok{panel.grid.major =}\NormalTok{ ggplot2}\SpecialCharTok{::}\FunctionTok{element\_line}\NormalTok{(}\AttributeTok{color =} \StringTok{"grey90"}\NormalTok{),}
    \AttributeTok{panel.grid.minor =}\NormalTok{ ggplot2}\SpecialCharTok{::}\FunctionTok{element\_line}\NormalTok{(}\AttributeTok{color =} \StringTok{"grey95"}\NormalTok{)}
\NormalTok{  )}

\FunctionTok{print}\NormalTok{(custom\_plot)}
\end{Highlighting}
\end{Shaded}

\subsection{Interpreting Output}\label{interpreting-output-2}

This workflow demonstrates how \texttt{tipr} can be part of a larger
data analysis pipeline. The output is a highly customized plot, clearly
showing the identified tipping point within the context of filtered and
smoothed data, with an enhanced aesthetic. This integration allows for
more complex data preparation and presentation.

\section{Advanced Features \&
Customization}\label{advanced-features-customization}

\subsection{Modifying Parameters}\label{modifying-parameters}

The \texttt{tipr::tip\_point()} function (and potentially others) might
have parameters that can be adjusted. For example, a \texttt{method}
parameter allows choosing different change-point detection algorithms.
Understanding these parameters and their impact is crucial for advanced
usage.

\begin{Shaded}
\begin{Highlighting}[]
\CommentTok{\# Example of method parameter:}
\CommentTok{\# Some change{-}point detection methods might be more suitable for different data types}
\CommentTok{\# or noise levels. Here, we demonstrate choosing \textquotesingle{}mean\textquotesingle{} vs \textquotesingle{}variance\textquotesingle{} based methods}
\CommentTok{\# if they were available within tipr\textquotesingle{}s hypothetical implementation.}

\CommentTok{\# Assume \textquotesingle{}tipr\textquotesingle{} could support different methods beyond just "changepoint"}
\CommentTok{\# For instance, methods focused on mean shifts vs. variance shifts.}
\CommentTok{\# As a placeholder, let\textquotesingle{}s illustrate how a different method *might* be used.}

\CommentTok{\# If tipr supported a method like "variance\_change"}
\CommentTok{\# variance\_tipping\_point \textless{}{-} data\_example |\textgreater{}}
\CommentTok{\#   tipr::tip\_point(}
\CommentTok{\#     data = \_,}
\CommentTok{\#     x = input\_variable,}
\CommentTok{\#     y = outcome\_variable,}
\CommentTok{\#     method = "variance\_change"}
\CommentTok{\#   )}
\CommentTok{\# print(variance\_tipping\_point)}

\CommentTok{\# Currently, \textquotesingle{}tipr\textquotesingle{} as described focuses on \textquotesingle{}changepoint\textquotesingle{}.}
\CommentTok{\# The main customization for \textasciigrave{}tip\_point\textasciigrave{} currently is the method.}
\CommentTok{\# For other functions, parameters might include \textasciigrave{}title\textasciigrave{}, \textasciigrave{}subtitle\textasciigrave{} for plots.}
\end{Highlighting}
\end{Shaded}

\subsection{Integration}\label{integration}

As seen in the advanced workflow, \texttt{tipr} naturally integrates
with \texttt{dplyr} for data manipulation and \texttt{ggplot2} for
visualization. This is a common pattern in R data analysis, where
specialized packages work together within the \texttt{tidyverse}
framework.

\begin{Shaded}
\begin{Highlighting}[]
\CommentTok{\# Integration with \textquotesingle{}tidyr\textquotesingle{} for data reshaping before tipr analysis}
\CommentTok{\# Imagine you have data in a wide format and need to convert it to long format}
\CommentTok{\# for tipr analysis across different groups.}

\NormalTok{wide\_data }\OtherTok{\textless{}{-}} \FunctionTok{data.frame}\NormalTok{(}
  \AttributeTok{time =} \DecValTok{1}\SpecialCharTok{:}\DecValTok{50}\NormalTok{,}
  \AttributeTok{group\_A\_value =} \FunctionTok{c}\NormalTok{(}\FunctionTok{rnorm}\NormalTok{(}\DecValTok{25}\NormalTok{, }\DecValTok{10}\NormalTok{, }\DecValTok{1}\NormalTok{), }\FunctionTok{rnorm}\NormalTok{(}\DecValTok{25}\NormalTok{, }\DecValTok{20}\NormalTok{, }\DecValTok{2}\NormalTok{)),}
  \AttributeTok{group\_B\_value =} \FunctionTok{c}\NormalTok{(}\FunctionTok{rnorm}\NormalTok{(}\DecValTok{30}\NormalTok{, }\DecValTok{5}\NormalTok{, }\FloatTok{0.5}\NormalTok{), }\FunctionTok{rnorm}\NormalTok{(}\DecValTok{20}\NormalTok{, }\DecValTok{15}\NormalTok{, }\DecValTok{1}\NormalTok{))}
\NormalTok{)}

\CommentTok{\# Reshape to long format for easier group{-}wise analysis with tipr}
\NormalTok{long\_data }\OtherTok{\textless{}{-}}\NormalTok{ wide\_data }\SpecialCharTok{|\textgreater{}}
\NormalTok{  tidyr}\SpecialCharTok{::}\FunctionTok{pivot\_longer}\NormalTok{(}
    \AttributeTok{cols =} \FunctionTok{c}\NormalTok{(group\_A\_value, group\_B\_value),}
    \AttributeTok{names\_to =} \StringTok{"group"}\NormalTok{,}
    \AttributeTok{values\_to =} \StringTok{"value"}
\NormalTok{  )}

\CommentTok{\# Now you could apply tipr::tip\_point() for each group using purrr or dplyr::group\_by()}
\CommentTok{\# For example, to find tipping points for each group:}
\CommentTok{\# (Requires group\_by and then applying the function, which is more advanced but shows integration)}

\CommentTok{\# tipping\_points\_by\_group \textless{}{-} long\_data |\textgreater{}}
\CommentTok{\#   dplyr::group\_by(group) |\textgreater{}}
\CommentTok{\#   dplyr::summarise(}
\CommentTok{\#     tipping\_point\_x = list(tipr::tip\_point(}
\CommentTok{\#       data = dplyr::cur\_data(),}
\CommentTok{\#       x = time,}
\CommentTok{\#       y = value,}
\CommentTok{\#       method = "changepoint")$tipping\_point)}
\CommentTok{\#   )}
\CommentTok{\# print(tipping\_points\_by\_group)}
\end{Highlighting}
\end{Shaded}

\subsection{Optimization}\label{optimization}

\begin{itemize}
\tightlist
\item
  \textbf{For large datasets:} \texttt{tipr}'s performance will depend
  on the underlying change-point detection algorithms. For very large
  datasets, consider downsampling or using methods optimized for big
  data if available. Pre-processing steps like smoothing (as shown in
  the advanced workflow) can also reduce noise and potentially improve
  algorithm efficiency.
\item
  \textbf{Vectorization:} R's strength lies in vectorized operations.
  Ensure your data preparation and any custom functions leverage
  vectorization instead of explicit loops where possible.
\item
  \textbf{Method Choice:} If \texttt{tipr} offered multiple
  \texttt{method} options for \texttt{tip\_point()}, some might be
  computationally more efficient than others. Choose the most
  appropriate and efficient method for your data.
\end{itemize}

\section{Troubleshooting and FAQs}\label{troubleshooting-and-faqs}

\subsection{Common Errors}\label{common-errors}

\begin{itemize}
\tightlist
\item
  \textbf{\texttt{Error:\ object\ \textquotesingle{}x\textquotesingle{}\ not\ found}
  or
  \texttt{Error:\ object\ \textquotesingle{}y\textquotesingle{}\ not\ found}}:
  This usually means you haven't correctly specified the \texttt{x} and
  \texttt{y} column names in \texttt{tipr::tip\_point()} or
  \texttt{tipr::plot\_tipping\_point()}. Ensure they match your data
  frame's column names exactly (case-sensitive).
\item
  \textbf{\texttt{Error\ in\ …\ :\ no\ valid\ changepoint\ found}}: This
  error can occur if the \texttt{method\ =\ "changepoint"} algorithm
  cannot detect a statistically significant change point in your data.
  This might happen if:

  \begin{itemize}
  \tightlist
  \item
    There is no actual abrupt shift in your data.
  \item
    The noise level is too high, masking the change.
  \item
    The data is too short or sparse for the algorithm to work
    effectively.
  \item
    Try visualizing your data first to see if a visual tipping point
    exists.
  \end{itemize}
\item
  \textbf{Plotting issues (e.g., plot not showing)}: Ensure you are
  \texttt{print()}ing the \texttt{ggplot2} object if you are running it
  outside of an interactive R session (e.g., in a script). Also, check
  for missing data (\texttt{NA} values) in your \texttt{x} or \texttt{y}
  variables, as \texttt{ggplot2} functions might not handle them
  gracefully without explicit handling.
\end{itemize}

\subsection{FAQs}\label{faqs}

\begin{itemize}
\tightlist
\item
  \textbf{Can ``tipr'' handle multiple tipping points?}

  \begin{itemize}
  \tightlist
  \item
    The primary \texttt{tipr::tip\_point()} function is designed to find
    the \emph{most significant} single tipping point. To find multiple
    points, you might need to apply the function iteratively on subsets
    of your data (as conceptually shown in the intermediate workflow) or
    use more advanced change-point detection packages that explicitly
    support multiple change points.
  \end{itemize}
\item
  \textbf{What if my data has a non-linear trend?}

  \begin{itemize}
  \tightlist
  \item
    \texttt{tipr::tip\_point()}'s \texttt{changepoint} method typically
    looks for shifts in underlying statistical properties, which can
    work even with non-linear trends if there's an abrupt change in the
    \emph{nature} of the trend. However, if the ``tipping point'' is a
    continuous acceleration or deceleration rather than a sharp shift,
    other modeling techniques might be more appropriate.
  \end{itemize}
\item
  \textbf{How do I choose the best \texttt{method} for
  \texttt{tipr::tip\_point()}?}

  \begin{itemize}
  \tightlist
  \item
    Currently, the primary method showcased is \texttt{"changepoint"}.
    If \texttt{tipr} expands to include more methods, the best choice
    depends on the specific characteristics of your data (e.g., whether
    you expect changes in mean, variance, or trend) and the underlying
    assumptions of each method. It's often best to experiment and
    compare results from different methods, guided by the theoretical
    basis of each.
  \end{itemize}
\end{itemize}

\section{Best Practices}\label{best-practices}

\begin{itemize}
\tightlist
\item
  \textbf{Visualize your data first:} Before using
  \texttt{tipr::tip\_point()}, always plot your data. This helps you
  visually inspect for potential tipping points and understand the
  general trend, which can guide your analysis and interpretation.
\item
  \textbf{Understand your data's context:} A statistically significant
  tipping point doesn't always imply practical significance. Understand
  the real-world context of your data to interpret the identified
  tipping points meaningfully.
\item
  \textbf{Clean and pre-process your data:} Noise, outliers, and missing
  values can significantly affect the accuracy of tipping point
  detection. Perform appropriate data cleaning, smoothing, or imputation
  steps before analysis.
\item
  \textbf{Iterate and validate:} Tipping point analysis is often
  iterative. Experiment with different parameters (if available) or
  methods, and consider validating your findings using out-of-sample
  data or domain expertise.
\item
  \textbf{Document your code:} Use clear comments to explain your steps,
  especially for complex workflows or custom modifications. This ensures
  reproducibility and understanding for yourself and others.
\item
  \textbf{Reproducible R code:} Always use \texttt{set.seed()} when
  generating random data to ensure your examples and analyses are
  reproducible. Use clear variable names and consistent code formatting.
\end{itemize}

\section{Conclusion}\label{conclusion}

The ``tipr'' package provides a valuable tool for identifying critical
\textbf{tipping points} in your data, allowing for proactive
decision-making and deeper insights into system dynamics. By
understanding and utilizing \texttt{tipr::tip\_point()},
\texttt{tipr::plot\_tipping\_point()}, and (conceptually)
\texttt{tipr::sensitivity\_analysis()}, intermediate R users can
effectively uncover significant thresholds in various applications.

\subsection{Key Takeaways}\label{key-takeaways}

\begin{itemize}
\tightlist
\item
  \textbf{Tipping points} represent critical thresholds where small
  changes lead to significant outcomes.
\item
  \texttt{tipr::tip\_point()} identifies these points using statistical
  change-point detection methods.
\item
  \texttt{tipr::plot\_tipping\_point()} provides clear visualizations of
  detected tipping points.
\item
  Sensitivity analysis is crucial for understanding the robustness of
  identified tipping points.
\item
  ``tipr'' integrates well with the \texttt{tidyverse} for comprehensive
  data analysis workflows.
\end{itemize}

\subsection{Next Steps}\label{next-steps}

\begin{itemize}
\tightlist
\item
  \textbf{Explore advanced change-point detection packages:} If
  \texttt{tipr}'s functionalities are insufficient for your specific
  needs, explore other R packages dedicated to change-point detection
  (e.g., \texttt{changepoint}, \texttt{bcp}) which might offer more
  methods or flexibility.
\item
  \textbf{Apply to your own datasets:} The best way to learn is by
  doing. Apply ``tipr'' to your own real-world datasets to identify
  tipping points relevant to your field.
\item
  \textbf{Deep dive into change-point theory:} For a more profound
  understanding, research the statistical and mathematical theories
  behind change-point detection algorithms.
\end{itemize}

\section{References and Resources}\label{references-and-resources}

\begin{itemize}
\tightlist
\item
  \textbf{Official Documentation:} For the most up-to-date and detailed
  information, always refer to the official ``tipr'' package
  documentation, typically available on CRAN or GitHub. (As this is a
  hypothetical package, please note that for a real package, you would
  look for its CRAN page or GitHub repository for official documentation
  and vignettes.)
\item
  \textbf{Helpful Learning Resources:}

  \begin{itemize}
  \tightlist
  \item
    \textbf{``An Introduction to Statistical Learning''} (with
    applications in R): While not specific to \texttt{tipr}, this book
    covers many foundational statistical concepts relevant to data
    analysis.
  \item
    \textbf{Online courses on R and Data Science:} Platforms like
    Coursera, edX, or DataCamp often have courses covering statistical
    modeling and data analysis in R.
  \end{itemize}
\item
  \textbf{Support Channels:}

  \begin{itemize}
  \tightlist
  \item
    \textbf{Stack Overflow:} A great place to ask specific R programming
    questions and troubleshoot errors.
  \item
    \textbf{RStudio Community:} A forum for R users to discuss and get
    help.
  \item
    \textbf{GitHub Issues:} If ``tipr'' were a real package, its GitHub
    repository would likely have an ``Issues'' section for reporting
    bugs or requesting features.
  \end{itemize}
\end{itemize}




\end{document}
